\section{Zadanie 2}
Tasujemy talię rozróżnialnych kart. W talii mamy $B$ białych kart i $C$ czarnych kart. Niech $r$ będzie ustaloną liczbą, taką że $1 \leq r \leq C$. Wykładamy karty dopóki nie odkryjemy $r$ czarnych kart. Niech $X$ będzie liczbą odkrytych białych kart. Jaki rozkład ma $X$?

\subsection*{Rozwiązanie}
\[
P(X=k) = \frac{\binom{k+r-1}{k}\cdot \binom{B+C-k-r}{B-k}}{\binom{B+C}{B}}
\]\\
Wyjaśnienie: \\
Rozważamy na ile sposobów mogą być ułożone wszystkie karty (nie przestajemy odkrywać kart po odkryciu $r$-tej czarnej)\\
Mamy $\binom{B+C}{B}$ możliwości na wybranie miejsc dla białych kart.\\
Teraz policzymy ułożenia kart, gdzie przez $r$-tą czarną jest $k$ białych. Na $\binom{k+r-1}{k}$ sposobów wybieramy miejsca na białe karty, które mają się znaleźć przed $r$-tą czarną, a na $\binom{B+C-k-r}{B-k}$ sposobów wybieramy miejsca na pozostałe białe karty.

\subsection*{Odpowiedź}
\[
\boxed{\text{$P(X=k) = \frac{\binom{k+r-1}{k} \cdot \binom{B+C-k-r}{B-k}}{\binom{B+C}{B}}$}}
\]
